\section{Supplementary Material}

This appendix records reproducibility details for the \benchmark{} benchmark
paper: training settings, regeneration commands, external sanity checks, and
qualitative probe examples.

\subsection{Training Configuration}

\begin{table}[ht]
\centering
\footnotesize
\begin{tabular}{ll}
\toprule
Setting & Value \\
\midrule
Backbone & Qwen2.5-VL-7B-Instruct \\
Training objective & Direct Preference Optimization \\
Adapter & LoRA \\
Epochs & 1 \\
DPO beta & 0.1 \\
LoRA rank & 16 \\
\answerdpo{} rows & 6{,}000 answer preference rows \\
\cepodual{} rows & 6{,}000 answer + 2{,}000 verifier rows \\
Distributed setup & ZeRO-2 \\
\bottomrule
\end{tabular}
\caption{Shared training setup for the two preference-tuned groups.}
\label{tab:app_train}
\end{table}

The completed local artifacts are:
\begin{itemize}
    \item canonical \dataset{} JSONL: 6{,}000 rows.
    \item \cepodual{} DPO JSONL: 8{,}000 rows.
    \item supported-evidence probe JSONL: 600 rows.
    \item wrong-evidence probe JSONL: 400 rows.
\end{itemize}

\subsection{Regeneration Commands}

Lightweight regeneration runs the CEPO data checker, the audit-sheet
completion script, the CEPO-Dual exporter, and the LLaMA-Factory data-prep
wrapper. Exact commands are listed in the repository README files.
Heavy training and inference should be submitted through Slurm:
\begin{quote}
\footnotesize
\texttt{submit\_cepo\_dual\_pipeline.sh}
\end{quote}
The final numbers in the paper were refreshed with:
\begin{quote}
\footnotesize
\texttt{score\_object\_eval.py}\\
\texttt{score\_cepo\_evidence\_probe.py}
\end{quote}

\subsection{External Results}

\begin{table}[ht]
\centering
\footnotesize
\begin{tabular}{lccc}
\toprule
Eval & Base & \answerdpo{} & \cepodual{} \\
\midrule
POPE random & 88.7 & 89.6 & \textbf{89.8} \\
POPE popular & 87.7 & 88.5 & \textbf{88.7} \\
POPE adversarial & 86.7 & 87.1 & \textbf{87.3} \\
AMBER discr. & 87.6 & \textbf{88.2} & \textbf{88.2} \\
\bottomrule
\end{tabular}
\caption{External yes/no sanity checks. Values are accuracy percentages under
the same deterministic parser.}
\label{tab:app_external}
\end{table}

\subsection{Qualitative Probe Examples}

\begin{table*}[ht]
\centering
\scriptsize
\begin{tabular}{@{}p{0.14\linewidth}p{0.34\linewidth}p{0.16\linewidth}p{0.13\linewidth}p{0.10\linewidth}@{}}
\toprule
Type & Candidate claim and evidence & Target & \cepodual{} output & Reading \\
\midrule
Supported attr. &
claim: vehicle is white; evidence: label=vehicle, attribute=white &
supported & supported & correct \\
Supported material &
claim: window is glass; evidence: label=window, attribute=glass &
supported & supported & correct \\
Wrong object &
claim: sink visible; evidence: label=potted plant &
wrong evidence & wrong evidence & correct \\
Wrong relation &
claim: glass left\_of grape; evidence swaps subject/object &
wrong evidence & contradicted & correct rejection \\
Wrong relation &
claim: man right\_of parking meter; evidence swaps subject/object &
wrong evidence & supported & relation failure \\
\bottomrule
\end{tabular}
\caption{Representative probe outputs. Examples are anonymized to local image
references and show why relation reversals remain the hardest slice.}
\label{tab:app_qual}
\end{table*}

\subsection{Scoring Notes}

Short-answer metrics use a deterministic yes/no parser. Outputs with explicit
uncertainty such as ``cannot determine'' or ``not sure'' are refusals; outputs
without a parseable yes/no decision are counted as other/invalid. The
evidence-probe scorer first tries to parse JSON and then falls back to labeled
fields. For wrong-evidence rows, both \texttt{wrong\_evidence} and
\texttt{contradicted} count as rejecting the candidate evidence, because both
labels indicate that the evidence should not be accepted as support for the
claim.
